% Transcription — fonds Alexandre Grothendieck, shelfmark 19, batch 5 (pages 81-95).
% Established from the facsimile archives/batches/19/batch-05.pdf (a mirror of
% https://grothendieck.umontpellier.fr/19.pdf), per the skill
% transcribe-grothendieck.
%
% The last batch of the folder: fifteen pages, not twenty, because the folder
% ends at 95.
%
% Pages 85 and 93 are blank versos — absent. What shows on them is the facing
% sheet (86 and 94 respectively) coming through the paper, not writing of their
% own.
%
% Page 81 carries six lines only: they finish the sentence that broke off at
% the foot of page 80, closing the "Questions sur les topos" of batch 4.
%
% Page 82 is the wrapper of the last piece and names it: "Catégories fibrées".
% Pages 83 to 95 are that piece, in three parts, each with a heading of the
% author's own: "Sites fibrés" (83-84), a boxed "Catégories Gamma et Lim"
% (86-92), and a boxed "Dualité pour catégories fibrées et cofibrées" (94-95).
%
% The middle part is the substantial one. Pages 86 to 88 set the square
% Lim E/B -> Gamma(E/B) against its base change and derive, by two Hom
% computations, the two adjoints F_*^c and F_+^c of F_c^*; the page 87
% proposition constructs the left adjoint by a fibrewise Lim-> ; page 88 draws
% the two corollaries on fibre functors and on generators. A notation is worth
% naming because the page depends on it: the author writes (-)_* for a LEFT
% adjoint and (-)_+ for a RIGHT one, and the two occur side by side.
% Pages 89 to 92 continue in pencil, on generators and cogenerators of Lim E/B
% and on when it is a U-topos; they are heavily overwritten and interlined, and
% are transcribed with many gaps. The bottom third of page 87 is cancelled by a
% long diagonal stroke and is largely unreadable.
%
% The dating is the inventory's own, for the whole shelfmark; nothing on these
% pages carries a date.
%
% Pass: Opus 5 (claude-opus-5), 2026-08-22 — first pass, unchecked against the pages by a human.
\documentclass[11pt,a4paper]{article}
% Le préambule commun à toutes les transcriptions du fonds Grothendieck.
%
% Il tient en une idée : **l'incertitude se compose, elle ne se résout pas.**
% Une transcription qui lisse un mot douteux a détruit la seule chose pour
% laquelle on la faisait — un lecteur doit pouvoir savoir, à chaque endroit,
% ce qui était sur le feuillet et ce qui a été ajouté. Les six macros qui
% suivent sont donc l'appareil critique, et non de la décoration.
%
% Les mêmes commandes sont reconnues par `scripts/render.mjs`, qui produit la
% vue de lecture du site. Une macro ajoutée ici doit l'être aussi là-bas,
% faute de quoi le rendu échouera bruyamment — ce qui est voulu : un rendu
% silencieusement faux est pire qu'un rendu absent.

\ProvidesPackage{grothendieck}[2026/08/08 Transcriptions du fonds Grothendieck]

\usepackage{amsmath,amssymb,amsthm}
\usepackage{mathrsfs}
\usepackage{stmaryrd}
% Les diagrammes commutatifs sont partout dans ce fonds : c'est la langue
% même de la géométrie algébrique de Grothendieck, et une transcription qui
% les rendrait en prose ne transcrirait rien.
\usepackage{tikz-cd}
% Français par défaut — c'est la langue des feuillets — mais l'anglais est
% chargé aussi : la traduction et le résumé sont des documents anglais, et la
% typographie française (espaces avant les deux-points) y serait une faute.
% Ces documents basculent par \selectlanguage{english} en tête de corps.
\usepackage[english,french]{babel}
\usepackage{fontspec}
\usepackage{geometry}
\geometry{a4paper,margin=25mm,marginparwidth=28mm}
\usepackage{marginnote}
\usepackage{xcolor}
% ulem, not soul. `soul`'s \st and \ul are built on a character-by-character
% scan that XeLaTeX's font machinery breaks: the document fails with an
% undefined control sequence rather than degrading. ulem does the same two jobs
% and is Unicode-engine safe. `normalem` stops it hijacking \emph, which the
% transcriptions use for Grothendieck's own emphasis.
\usepackage[normalem]{ulem}
\usepackage{hyperref}
\hypersetup{colorlinks=true,linkcolor=black,urlcolor=black,pdfborder={0 0 0}}

% --- Métadonnées du lot -----------------------------------------------------
% Reprises telles quelles par le script de rendu, qui les affiche en tête de
% la vue de lecture. Elles fixent ce que le fichier prétend couvrir : sans
% elles, une transcription flotte sans point d'attache dans un fonds de seize
% mille feuillets.

\newcommand{\folder}[1]{\gdef\grothFolder{#1}}
\newcommand{\batch}[1]{\gdef\grothBatch{#1}}
\newcommand{\leaves}[2]{\gdef\grothFirst{#1}\gdef\grothLast{#2}}
\newcommand{\foldertitle}[1]{\gdef\grothFoldertitle{#1}}
\gdef\grothFolder{?}\gdef\grothBatch{?}\gdef\grothFirst{?}\gdef\grothLast{?}\gdef\grothFoldertitle{}

% --- Le résumé --------------------------------------------------------------
%
% Il ouvre la lecture modernisée, sous le titre « Résumé », et il est composé
% en petit. Ce n'est pas un ornement : c'est le seul passage du document qui
% ne soit pas des mathématiques, et un lecteur doit pouvoir le reconnaître
% avant de l'avoir lu — puis le sauter s'il connaît déjà le sujet, ou s'y
% arrêter s'il ne le connaît pas. Le filet qui le suit marque la reprise du
% corps.

\newenvironment{resume}
  {\small\setlength{\parskip}{0.45em}}
  {\par\medskip\hrule\medskip}

% --- Mots-clés ---------------------------------------------------------------
%
% En anglais, en clôture du résumé de la lecture modernisée : le vocabulaire
% moderne sous lequel chercher ce que les feuillets construisent. Le manifeste
% du site (`npm run manifest`) les extrait pour étiqueter la cote entière —
% cette ligne est la source unique des tags, il n'y a pas de fichier à côté.
\newcommand{\keywords}[1]{\par\smallskip\noindent
  {\footnotesize\textsc{Keywords} --- \textit{#1}}\par}

% --- L'appareil critique ----------------------------------------------------

% Le numéro de feuillet, en marge. C'est l'ancre qui permet de régler un
% désaccord sur une formule en regardant *un* feuillet plutôt qu'une chemise
% de six cents. Le site s'en sert aussi pour tourner les pages du fac-similé
% quand on fait défiler la transcription.
\newcommand{\leaf}[1]{%
  \marginnote{\footnotesize\color{gray!70}\textsf{#1}}%
  \label{leaf:#1}\ignorespaces}

% Illisible : on ne devine pas. Un mot inventé qui se lit comme les autres est
% le pire résultat possible.
\newcommand{\ill}{\textcolor{red!60!black}{[\,\ldots\,]}}

% Lecture douteuse : le mot est donné, et signalé comme incertain.
\newcommand{\uncertain}[1]{\uline{#1}}

% Ajout de l'éditeur — une lettre manquante, un mot sous-entendu.
\newcommand{\add}[1]{\textcolor{blue!50!black}{[#1]}}

% Biffé par Grothendieck lui-même. Ce qu'il a rayé fait partie du document :
% une rature dit souvent où la pensée a changé de direction.
\newcommand{\struck}[1]{\sout{#1}}

% Note du transcripteur, sur l'état matériel du feuillet.
\newcommand{\note}[1]{\par\smallskip{\small\color{gray!60!black}\textit{[#1]}}\par\smallskip}

% Note marginale de Grothendieck, distincte du corps du texte.
\newcommand{\marginal}[1]{\par\smallskip\noindent\hspace{1em}%
  {\small\color{gray!60!black}#1}\par\smallskip}

% --- Titre ------------------------------------------------------------------

\AtBeginDocument{%
  \begin{center}
    {\large\bfseries Fonds Alexandre Grothendieck --- cote n\textsuperscript{o}~\grothFolder}\\[2pt]
    {\itshape\grothFoldertitle}\\[4pt]
    {\small feuillets \grothFirst--\grothLast\ (lot \grothBatch)}\\[2pt]
    {\footnotesize\color{gray!60!black}Transcription établie d'après le fac-similé numérisé
     par l'Université de Montpellier.\\
     Les lectures incertaines sont soulignées, les illisibles notées [\,\ldots\,],
     les ajouts entre crochets.}
  \end{center}
  \vspace{1em}}

\endinput


\folder{19}
\batch{5}
\pages{81}{95}
\dating{[à partir de 1958-1973]}
\watermark{Édition de démonstration}
\foldertitle{Topos : notes manuscrites (s.d.)}

\begin{document}

\section*{Questions sur les topos (fin, page 81)}

\page{81}\note{les six lignes de cette page achèvent la phrase interrompue au
bas de la page 80}

\uncertain{types} \ill{} formés des couples d'un ens.\ $A$ \struck{\ill{}} muni
d'une \uncertain{involution} bijectif $w$, et d'une rétraction de $A/w$ sur une
partie de $A^{w} \subset A/w$. Le foncteur $f$ est le foncteur qui associe à un
\ill{} le \ill{} ($E$ muni de $\mathrm{id}_{E}$, $E$, $\mathrm{id}_{E}$)

\section*{Catégories fibrées (pages 82 à 95)}

\page{82}
Titre porté sur la chemise : « Catégories fibrées ». Un pointage d'archiviste,
« 22 p », est écrit au crayon en regard.\note{les feuillets qui suivent, 83 à
95, n'en comptent que treize, dont deux versos vierges}

\subsection*{Sites fibrés}

\page{83}

$E \to B$ catégorie fibrée, avec les \struck{foncteurs} catégories fibres des
sites, les \struck{foncteurs} $f^{*}$ des morphismes de sites [conditions
suffisantes pour \ill{} autres choses]. On met sur $E$ la top.\ la moins fine de
celles qui rend continus les foncteurs $E_{b} \to E$. Les familles couvrantes
sont décrites par IV 3.2.

Supposons $B$, $E \in \mathcal{U}$, \struck{\ill{}} (que signifie ?) les top.\
sur $E_{b}$ \uncertain{Abondant}. Plongeons \ill{} $E_{b}$ en
$\widetilde{E}_{b}$, i.e.\ trouver une nouvelle catégorie fibrée
$\overline{E}$ : fibres les $\widetilde{E}_{b}$ ; un f.\ \struck{\ill{}}, un
morphisme de $F \in \overline{E}_{x}$ et $G \in \widetilde{E}_{y}$ est un
morphisme dans $\widetilde{E}_{x}$

\[ F \longrightarrow f^{*}(G) \]

$x \longrightarrow y$ [l'on désignant aussi par $f^{*}$ l'ext.\ de $f^{*}$ aux
$\widetilde{E}_{b}$].

Déterminons le type $\widetilde{E}$, dans un \ill{} \ill{}
\struck{$\widetilde{E} \simeq \varprojlim_{X} (E_{p(X)})_{/X} = \varprojlim_{x
\in \mathrm{Ob}\,B} E_{x}$}\note{la formule est encadrée, puis barrée} mais
étant entendu que $E$ est la catégorie fibrée dont \ill{} aux sections sur $X
\in \mathrm{Ob}\,E$, \ill{} la fibre, \ill{} associé au « point \ill{} » défini
par $X$ en $E$, savoir

\page{84}

$(\widetilde{E}_{p(x)})_{/X}$ \struck{$\simeq \widetilde{E}_{x}$},
\struck{\ill{}} qui a se \ill{} cartésienne quand $X$ parcourt une \ill{}
(« \ill{} $\geq x$ ») $E_{x}$. De façon précise

\[ \widetilde{E}^{\circ} \simeq
   \underline{\Gamma}^{\,\text{vert-cart}}\bigl(B,\;
   X \mapsto (\overline{E})^{\circ}_{p(x)/X}\bigr) \]

où l'exposant « vert-cart » signifie qu'on se borne aux sections qui sont
cartésiennes au-dessus des flèches verticales. Utilisant une formule
d'associativité évidente, le deuxième membre s'interprète comme

\[ \underline{\Gamma}\Bigl(B,\; x \mapsto
   \underline{\Gamma}^{\,\mathrm{cart}}\bigl(E_{x},\;
   X \mapsto \underbrace{(\overline{E}_{x})^{\circ}_{/X}}
   _{\textstyle \widetilde{E}^{\circ}_{x} = (\overline{E})^{\circ}_{x}}
   \bigr)\Bigr) \]

Donc

\[ \widetilde{E} \simeq \underline{\Gamma}(B, \overline{E}^{\Delta})^{\circ} \]

où $\overline{E}^{\Delta}$ désigne la catégorie fibrée déduite de
$\overline{E}$ en prenant \struck{\ill{}} \uncertain{groupoïdes} les fibres de
$\overline{E}$.

L'inclusion $\varprojlim\bigl(E,\; X \mapsto
(\overline{E})_{p(x)/X}\bigr) \to \underline{\Gamma}^{\,\mathrm{vert}}
(\widetilde{E}, \ldots)$ vient alors \ill{} que l'inclusion

\[ \underline{\Gamma}\,\varprojlim(B, \overline{E}^{\Delta})^{\circ}
   \simeq \varprojlim(B, \overline{E})^{\Delta}
   \subset \underline{\Gamma}(B, \overline{E})^{\Delta} \]

\subsection*{Catégories $\underline{\Gamma}$ et $\varprojlim$}

\page{86}\note{le titre est encadré en haut de la page}

$E \to B$ catégorie fibrée.

\begin{tikzcd}
& \varprojlim E/B \arrow[r, hook, "i"] \arrow[d, "F^{*}_{c}"']
  & \underline{\Gamma}(E/B) \arrow[d, "F^{*}"] \\
X' \in B & \varprojlim E'/B' \arrow[r, hook, "i'"]
  & \underline{\Gamma}(E'/B')
\end{tikzcd}

\begin{align}
\mathrm{Hom}\bigl(X', F^{*}_{c}(Y)\bigr)
  &= \mathrm{Hom}\bigl(i'(X'), F^{*}(i(Y))\bigr) \notag \\
  &= \mathrm{Hom}\bigl(F_{*}\,i'(X'), i(Y)\bigr) \notag \\
  &= \mathrm{Hom}\bigl(\underbrace{i_{*}F_{*}\,i'(X')}
     _{\textstyle F^{c}_{*}(X')}, Y\bigr) \notag
\end{align}

\textbf{①} Si $\exists$ adjoint à g.\ $i_{*}$ de $i$, et si $\exists$ $F_{*}$
adjoint à g.\ de $F^{*}$ (p.ex.\ si $E$ est \emph{assez} cofibrée sur $B$, et
les « petites » $\varprojlim$ existent dans les fibres) alors $\exists$ adjoint
à gauche $F^{c}_{*}$ de $F^{*}_{c}$.

\begin{align}
\mathrm{Hom}\bigl(F^{*}_{c}(Y), X'\bigr)
  &= \mathrm{Hom}\bigl(F^{*}(i(Y)), i'(X')\bigr) \notag \\
  &= \mathrm{Hom}\bigl(i(Y), F_{+}\,i'(X')\bigr) \notag \\
  &= \mathrm{Hom}\bigl(Y, \underbrace{i_{+}F_{+}\,i'(X')}
     _{\textstyle F^{c}_{+}}\bigr) \notag
\end{align}

\note{l'auteur écrit ici l'indice « $+$ » là où le calcul précédent portait un
« $*$ » : dans toute cette page $(-)_{*}$ désigne un adjoint \emph{à gauche} et
$(-)_{+}$ un adjoint \emph{à droite}}

\textbf{②} Si $\exists$ adjoint à dr.\ $i_{+}$ \struck{du} de $i$, et si
$\exists$ adj.\ à dr.\ $f_{+}$ de $F^{*}$ (p.ex.\ si les petites $\varprojlim$
existent en

\page{87}

les fibres de $E/B$), alors $\exists$ adjoint à droite $F^{c}_{+}$ de
$F^{*}_{c}$.

\textbf{③ Prop.} Supposons que pour tout $b \in \mathrm{Ob}\,B$, la limite
inductive $\varinjlim_{c \in \mathrm{Ob}\,{}^{b}\!/B} f^{*}(X(c))$ dans $E_{b}$
existe, et supposons que cette $\varinjlim$ commute à tous les changements de
base $u^{*}$ $(a : b' \to b)$. Alors le foncteur adjoint à gauche
$\dot{F}^{\varepsilon}_{*}$ de $E$ : $\varprojlim E/B \longrightarrow
\underline{\Gamma}(B/E)$ existe\note{$\underline{\Gamma}(B/E)$ est bien ce qui
est écrit ; partout ailleurs l'auteur note $\underline{\Gamma}(E/B)$}, et est
donné par

\[ \dot{F}^{\varepsilon}_{*}(X)(b)
   = \varinjlim_{x \in \mathrm{Ob}\,{}^{b}\!/B} f^{*}(X(x)) \]

[Le morphisme d'adjonction

\[ \dot{F}^{\varepsilon}_{*}\,\dot{F}^{\varepsilon}_{*}(X) \longrightarrow X \]

\ill{} les vérifications évidentes]

\ill{} si de plus $E \to X$ est cofibrant, \ill{} $b^{*}_{c} : \varprojlim E/B
\to E_{b}$ \ill{} des \uncertain{adjoints} \ill{}\note{le dernier tiers de la
page — une dizaine de lignes, plus une note marginale portant l'encadré
$(\varprojlim E/B)^{\circ} \to \underline{\Gamma}(E/B)^{\circ}$ et un
« N.B. » — est barré d'un long trait diagonal et reste illisible}

\ill{} à droite des $b^{*}_{c} : \varprojlim E/B \to E_{b}$ existant également
et se construisent explicitement, s'il en va ainsi des $f^{*}$.

\page{88}

\textbf{Cor.} Si de plus les $b_{*}$ existent $(b \in \mathrm{Ob}\,B)$ [p.ex.\
si $p$ est aussi cofibrant, et les « petites » sommes dans les $E_{b}$
existent), alors les \struck{\ill{}} « foncteurs fibres » $\varprojlim E/B
\xrightarrow{\;b^{*}_{c}\;} E_{b}$ ont un adj.\ à gauche $b^{c}_{*}$.

\textbf{Cor.} Supposons que $\underline{\Gamma}$ ait une famille de
générateurs \add{stricts} $(X_{\lambda})$, alors les $i(X_{\lambda})$ forment
une famille de générateurs (stricts) de $\varprojlim$. En particulier, si $p$
cofibrant, et si les $E_{b}$ ont des \add{petites} sommes et de petites
familles de générateurs stricts, \struck{il en est \ill{} de $\underline{\Gamma}$}),
alors ces conditions sont vérifiées pour $\varprojlim$, et on trouve un syst.\
de générateurs (stricts) (en prenant pour la \ill{} $E_{b}$ des gén.\ stricts
de $E_{b}$, $b \in \mathrm{Ob}\,B$, et prenant

\[ \bigcup_{b \in \mathrm{Ob}\,B} b^{c}_{*}(E_{b}). \]

\page{89}

[ Si $\varinjlim$ \add{petites} \ill{} de $E_{b}$, et les $f^{*}$ y commutent :
$\varprojlim(E/B)^{\circ}$ \ill{} adjoint \ill{}

Si petites $\varprojlim$ \ill{} dans $E_{b}$, et les $f^{*}$ y commutent :

\[ \varprojlim{}^{\circ}(E^{\circ}/B) \xrightarrow{\;d^{\circ}\;}
   \underline{\Gamma}(E^{\circ}/B) \]

a un adjoint à gauche, donc

\[ \varprojlim(E/B) \xrightarrow{\;j\;}
   \underline{\Gamma}(E^{\circ}/B)^{\circ} \]

a un adjoint à droite $j_{+}$.

Si on demande des générateurs de $L$, il y a lieu d'utiliser $i_{*}$, car pour
des générateurs $X_{\lambda}$ de $\underline{\Gamma}(E/B)$, comme

\[ \mathrm{Hom}_{\underline{\Gamma}}\bigl(X_{\lambda}, i(Y)\bigr)
   \simeq \mathrm{Hom}_{L}\bigl(i_{*}(X_{\lambda}), Y\bigr), \]

\ill{} celles des $i_{*}(X_{\lambda})$ \ill{} dans $L$. Or, il n'y a pas de
\uncertain{relation} \ill{} dans $\underline{\Gamma}$, le cor.\ \ill{} \ill{}
si \ill{} $\varinjlim$ \ill{} $B^{\circ}$ des \ill{} dit qu'il existe \ill{}
$p$ cofibrant et \ill{} $E_{b}$, et que \ill{} gén.\ petits ens.\ gén.\ \ill{}
existe dans les $E_{b}$, les $f^{*}$ \ill{} de générateurs dans les $E_{b}$

\page{90}

Si on demande des cogénérateurs de $L$, i.e.\ des générateurs de $L^{\circ}$,
il y a lieu d'utiliser $j^{\circ}$ \ill{} à $\varphi^{\Delta}$, et on trouve
comme condition suff.\ que $j^{\circ}_{*}$ existe et qu'il y ait
\uncertain{petit} syst.\ de gén.\ dans $\underline{\Gamma}(E^{\Delta}/B)$,
(p.ex.\ que les petites $\varprojlim$ existent dans les $E_{b}$, que les
foncteurs $f^{*}$ y commutent, que les $E_{b}$ aient des petits ens.\ de
cogénérateurs]

\page{91}

\textbf{§} \struck{\ill{}}

Si dans les catégories fibres, certaines \add{types de} $\varprojlim$ ou
$\varinjlim$ sont représentables, et si les $f^{*}$ y commutent, alors les
\uncertain{mêmes} types de $\varprojlim$ sont représentables dans $\varprojlim
E/B$, et les foncteurs fibres (plus \ill{} les foncteurs changements de base,
$B' \to B$) y commutent.

Ainsi, si les $E_{b}$ sont \uncertain{préadditives} (\ill{} additives,
\add{abéliennes}) \add{et si les $f^{*}$ \ill{} respectent ces struct.}, il en
est de même de $\varprojlim E/B$, si $E_{b}$ \ill{} : $\varprojlim$
filtrantes\ill{} $= \ldots \varprojlim(B)$, ex.\ \ill{} et si les $H_{i}$,
$\varprojlim$ \ill{} y sont représentables (\ill{}) \ill{} dans $\varprojlim
E/B$. De \ill{}, si \ill{} les $E_{b}$ \ill{} conditions a), b), c) de \ill{}
et les $f^{*}$ des \ill{} \ill{} $\mathcal{U}$-topos \ill{} $= \ldots$
$\varprojlim E/B$.

Pour \uncertain{pouvoir} conclure que si les $E_{b}$ sont des
$\mathcal{U}$-topos \add{et les $f^{*}$ des morphismes de tels}, il en est de
même de $\varprojlim E/B$, il reste à donner des conditions \ill{} que cette
catégorie \ill{} une \add{petite} famille de gén.\ \ill{}
$\mathcal{U}$-catégorie, et \ill{}. Le premier est \uncertain{immédiat} \ill{}
que $B$ \ill{} « ens.\ petits ». Pour le dernier, il en est \ill{}

\page{92}\note{feuillet au crayon, très surchargé d'interlignes ; la
transcription en est très lacunaire}

\ill{} à \uncertain{supposer} que les \ill{} soient \uncertain{cofibrées},
i.e.\ que les \uncertain{adjoints à gauche} des $f^{*}$ existent, [qu'il n'y a
pas lieu \ill{} de notre \ill{}, mais plutôt $f_{!}$] \add{i.e.\ les $f^{*}$
\ill{} associés à des compléments $f_{!}$}, (\ill{} sont \ill{} qu'il vient
\ill{} \add{petites} \ill{} que les $f^{*}$ commutent aux $\varprojlim$)
\ill{} qui \uncertain{exige} \ill{} que \ill{} \ill{} en calculant fibre par
fibre. \ill{} exact \ill{} \ill{} \ill{} $f^{*}$ \ill{} commutent \ill{} Donc
\ill{} \ill{} \add{« ens.\ petits »} général, \ill{} $B$ \ill{} \ill{} il \ill{}
cependant \ill{} \ill{} \ill{} et générateurs \ill{} qu'il y a \ill{}
\uncertain{détail} \ldots. Donc \ill{} à rédiger \ill{} $\mathcal{U}$-topos
\ill{} $\mathcal{U}$-ob-topos. \ill{} bien des $\varprojlim$, mais les \ill{}
$\mathcal{U}$ y a des \ill{} faisant \ill{} calcul \ill{} générale, \ill{} un
peu \uncertain{fastidieux}.

\subsection*{Dualité pour catégories fibrées et cofibrées}

\page{94}\note{le titre est encadré en haut de la page}

Catégorie $E \xrightarrow{\;p\;} B$ sur une autre.

1) La catégorie est \emph{fibrée} sur $p$ \uncertain{fibrant}, \ill{} \ill{} $f
: x \to y$ dans $B$, i.e.\ pour définir $f^{*} : E_{y} \to E_{x}$, avec
\uncertain{transitivité essentielle} :

\[ \mathrm{Hom}_{f}(X, Y) \simeq \mathrm{Hom}_{\mathrm{id}_{x}}
   \bigl(X, f^{*}(Y)\bigr) \]

\ill{} \ill{} \ill{} aux pseudo-foncteurs $F : B^{\circ} \to (\mathrm{Cat})$.

2) La catégorie est \emph{cofibrée}, ou $p$ cofibrante, si $p^{\circ}$ est
fibrante, i.e.\ si pour tout $f : x \to y$ dans $B$, existe un foncteur $f_{*}
: E_{x} \to E_{y}$, avec transitivité essentielle :

\[ \mathrm{Hom}_{f}(X, Y) \simeq \mathrm{Hom}\,\varphi_{p}(X', Y) \}. \]

\note{lecture douteuse du second membre ; la symétrie avec 1) demanderait
$\mathrm{Hom}_{\mathrm{id}_{y}}(f_{*}(X), Y)$}

Ils donnent \uncertain{naissance} \ill{} un pseudo-foncteur

\[ F : B \longrightarrow (\mathrm{Cat}) \]

\emph{Lien} : Si $p$ est fibrant, pour que (pour $f : x \to y$ donnée) $f_{*}$
existe, il f.\ et s.\ que $f^{*}$ ait un adjoint à g. Pour que $p$ soit
cofibrant, il f.\ et s.\ que les $f^{*}$ aient tous des adjoints à g.

b) Si $p$ est cofibrant, \ill{} que $f^{*}$ existe, il f.\ et s.\ pour que
$f_{*}$ ait un adjoint à droite. Pour que $p$ soit fibrant, il f.\ et s.\ que
\ill{} adjoints à droite.

La catégorie \emph{opposée relative} d'une catégorie fibrée $p : E \to B$ n'est
définie seulement à équivalence près, c'est une autre catégorie fibrée,
associée au pseudo-foncteur

\page{95}

\[ F^{\mathrm{opp}} : B \longrightarrow (\mathrm{Cat}) \]

défini par

\begin{align}
F^{\mathrm{opp}}(a) &= F(a)^{\circ} \notag \\
F^{\mathrm{opp}}(f) &= F(f)^{\circ} \notag
\end{align}

\ill{} formule \ill{} pour \ill{} les \ill{} de transitivité :

Il n'est pas évident qu'il en existe une description intrinsèque à \ill{} près,
en évitant \ill{} recours aux pseudo-foncteurs.

Pour que $(E/B)^{\mathrm{opp}}$ soit aussi cofibrée, il f.\ et s.\ que les
foncteurs $f^{*} : E_{y} \to E_{x}$ aient des adjoints à \emph{droite}
(condition \uncertain{apparaît} : celle pour que $E/B$ soit cofibrée).

À \emph{retenir} \uncertain{que} ce pseudo-foncteur définit deux catégories
(fibrées) bien distinctes sur $B$.

\end{document}
